\chapter{The expression language}
\label{sec:language}


The language has two syntactic layers.

The \emph{internal layer} represents every expression as a plain
Lisp S-expression.  This representation is time-tested and deliberately
simple: it requires no parser, supports straightforward recursive
processing, and is directly inspectable and constructible at the Scheme
REPL.  All prover machinery---inference rules, substitution, pattern
matching, macetes, arithmetic evaluation---operates on this layer.
Users are not required to engage with it directly; however, in the
tradition of early Lisp, they are free to do so.  Constructing or
manipulating S-expressions by hand, loading axioms as quoted lists,
or writing Scheme code that builds formulas programmatically are all
legitimate and supported uses.

The \emph{string syntax} is the surface notation accepted by
\texttt{make-wff-from-string} and emitted by \texttt{wff->string}
/ \texttt{pp}: binary operators are written infix, multi-argument
terms use functional notation, quantifiers use the binding-list form,
and set expressions use brace notation.
Section~\ref{sec:parser-printer} describes both layers in full.
The examples in this manual use string syntax for user-facing
formulas; internal S-expression forms appear only in theory and
structure definitions.

MIT~Scheme reads symbols case-insensitively (folding to lower case),
so \texttt{FORALL}, \texttt{Forall}, and \texttt{forall} are the same
symbol at the Scheme level.

As should be the case in any formal system, we distingish between
syntax and semantics. In what follows, we distinguish two distinct
types of syntactic categories, each with its own semantics: formulas
and terms. These syntactic categories can be viewed in either the
internal syntax or the string syntax. In the presentation below, we
adhere to the string syntax.


\section{Formulas}

Semantically, a formula is a syntactic expression that can have only
two values: TRUTH or FALSEHOOD. 
%
\begin{remark}\label{remark:27137-52766-301830}
%
Note that ``UNDEFINED'' is not a possible value for a
formula. However, the value of any particular formula may not be
known. In fact, since we are silent about how elements in our system
may be interpreted in a pure von Neumann Bernays theory, different
interpretations may yield different values.
%
\end{remark}

Formulas can be produced from other ones by basic logical operators,
including quantification over variables that range over objects.
%
In the following table, the symbols $p$, $q$ are not variables that
range over formulas, but rather placeholders that can instantiated
with well-formed formulas.

\begin{itemize}[label=\textendash, leftmargin=1cm]
\item \texttt{truth} propositional truth $\top$
\item \texttt{falsity} propositional falsity $\bot$
\item \texttt{not($p$)} negation $\lnot p$
\item \texttt{$p$ and $q$} conjunction $p \land q$
\item \texttt{$p$ or $q$} disjunction $p \lor q$
\item \texttt{$p$ implies $q$} implication $p \Rightarrow q$
\item \texttt{$p$ iff $q$} biconditional $p \Leftrightarrow q$
\item \texttt{forall([$b_1$, \ldots, $b_n$], $p$)} universal quantification; see Section~\ref{sec:multiquant}
\item \texttt{forsome([$b_1$, \ldots, $b_n$], $p$)} existential quantification; see Section~\ref{sec:multiquant}
\item \texttt{$a$ = $b$} strict equality: both defined and equal
\item \texttt{$a$ == $b$} quasi-equality: both defined and equal, or both undefined
\item \texttt{$a$ in $b$} membership $a \in b$
\item \texttt{$A$ subset $B$} subclass $A \subseteq B$; see Section~\ref{sec:class-set-discipline}
\end{itemize}

\noindent
The quantifiers \texttt{forall} and \texttt{forsome} range over \emph{all
classes} (the system is single-sorted).  There is no restriction to sets.

\subsection{Quantification}
\label{sec:multiquant}

Both \texttt{forall} and \texttt{forsome} require a bracket-enclosed
\emph{binding list}, even for a single variable:

\begin{itemize}[label = $\circ$, leftmargin=1cm]
\item \texttt{forall([binding, ...], body)}

\item \texttt{forsome([binding, ...], body)}
\end{itemize}

\noindent
\texttt{forall(x, x = x)} is a syntax error; the correct form is
\texttt{forall([x], x = x)}.

\noindent
Each \emph{binding} in the string syntax is one of:

\begin{itemize}[label = $\circ$, leftmargin=1cm]
\item \texttt{$x$} unrestricted variable ranging over all classes
\item \texttt{$x$ in $A$} variable $x$ restricted to class $A$
\item \texttt{[$n_1$, \ldots, $n_k$] in $A$} tuple destructuring: fresh $t \in A$, components bound to \texttt{nth(1,~$t$)}, \ldots, \texttt{nth($k$,~$t$)}
\end{itemize}

\noindent
In the \emph{S-expression syntax} a binding can additionally group
several unrestricted variables into a single spec:
\[
  \texttt{(FORALL ((}x_1\ x_2\ \ldots\ x_n\texttt{) ...) body)}
\]
The grouped binding \texttt{($x_1$ $x_2$ \ldots\ $x_n$)} is exactly
equivalent to the sequence of single-variable bindings
\texttt{($x_1$) ($x_2$) \ldots\ ($x_n$)}.
The utility functions \texttt{vnb-unwind} and \texttt{vnb-wind}
convert between the two forms:

\begin{VNBsnip}
(vnb-unwind '(FORALL ((x y z)) body))
;=> (FORALL ((x)) (FORALL ((y)) (FORALL ((z)) body)))

(vnb-wind '(FORALL ((x)) (FORALL ((y)) (FORALL ((z)) body))))
;=> (FORALL ((x) (y) (z)) body)
\end{VNBsnip}

\noindent
Bindings are processed left to right, outermost first.  Thus in the
\emph{internal syntax} they expand to fully nested single-variable form
as follows:

\begin{VNBsnip}
forall([x, y in NN, z], y >= 0)
; => (FORALL x (FORALL y (IMPLIES (IN y NN) (FORALL z (>= y 0)))))

forall([[E, +, *] in METRIC], IS-METRIC(E, +, *))
; => (FORALL t_0 (IMPLIES (IN t_0 METRIC)
;                         (IS-METRIC (NTH 1 t_0) (NTH 2 t_0) (NTH 3 t_0))))
\end{VNBsnip}

\paragraph{Display.}
The printer reconstructs binding-list form from the internal nested
representation.  Consecutive same-type quantifiers are collapsed into
a single binding list; a change of quantifier type stops the collapsing.
For example, \texttt{forall([x in NN], forall([y], p))} prints as
\texttt{forall([x in NN, y], p)}.

\paragraph{Variable shadowing.}
Bound variable names may be reused within nested quantifiers.
The inner binding shadows the outer; \texttt{fresh-var} in
\texttt{subst-free} ensures substitution never captures a free variable.
Such formulas are valid but confusing; users should avoid them.

\subsection{Equality and quasi-equality}

Terms may be \emph{undefined} (e.g.\ an application \texttt{f(x)} when
\texttt{f} is not a function applicable to \texttt{x}).  The system
provides two equality predicates:

\paragraph{Strict equality \texttt{$a$ = $b$}.}
True iff both \texttt{a} and \texttt{b} are defined and equal.
False if either is undefined.  In particular:
\begin{itemize}[label=$-$, leftmargin=1cm]
\item \texttt{george(s) = 75} is \emph{false} when \texttt{george}
  is not a function.
\item \texttt{george(s) = george(s)} is \emph{false} for the same
  reason.
\item A term $t$ is defined iff $t = t$:
  $\mathit{DEF}(t) \;\equiv\; (t = t)$.
\end{itemize}

\paragraph{Quasi-equality \texttt{$a$ == $b$}.}
True iff both \texttt{a} and \texttt{b} are defined and equal, \emph{or}
both are undefined.  Key properties:
\begin{itemize}[label=$-$, leftmargin=1cm]
\item \texttt{t == t} is \emph{always true}, for any term $t$,
  defined or not.
\item \texttt{george(s) == george(s)} is \emph{true}.
\item \texttt{$a$ = $b$} implies \texttt{$a$ == $b$}, but not vice versa.
\item \texttt{$a$ == $b$} and $\mathit{DEF}(a)$ together imply
  \texttt{$a$ = $b$}.
\end{itemize}

\noindent
Undefinedness propagates through all operations: if \texttt{george(s)}
is undefined then \texttt{george(s) + 1} is also undefined,
and \texttt{george(s) + 1 = 2} is false.

Variables and class constants (\texttt{ORD}, \texttt{RR},
\texttt{EMPTY-SET}, \ldots) are always defined.  The application
\texttt{f(x)} is defined when \texttt{f in fun(A, B)} and
\texttt{x in A} are provable for some $A$, $B$.

\paragraph{Reflexivity in proofs.}
\texttt{cmd-quasi-reflexivity} closes any goal of the form
\texttt{a == a} unconditionally.
\texttt{cmd-reflexivity} closes \texttt{a = a} only when \texttt{a} is
demonstrably defined.

\section{Class constructors}
\label{sec:class-set-discipline}

Proper VNB has two kinds of object: \emph{sets} (elements of
$\mathit{SET}$) and \emph{proper classes} (collections that can have
members but are not themselves members of anything).  We retain that
distinction and add a third kind, \emph{functoids}.

A \textbf{functoid} is an operation specified by a \emph{signature}
\[
  A_1 \times A_2 \times \cdots \times A_n \;\longrightarrow\; B
\]
where $A_1, \ldots, A_n$ are domain classes and $B$ is the codomain
class.  None of them need be sets.

\begin{remark}
The motivation comes from category theory.  A functor is conceptually a
mapping between collections of mathematical structures, but its domain
is typically a proper class (e.g.\ the class of all groups, or all
topological spaces).  In a reductionist treatment one could introduce a
distinguished class $\mathit{FUNCTOID} \subseteq \mathit{CLASS}$,
whose members are ordered pairs encoding signature and rule.  This
system instead treats functoids as a new kind of object, coordinate
with sets and classes rather than reducible to them.  The principal
consequence is that the formal language does not quantify over
functoids: ``all functoids $F$ satisfy\ldots'' is metalanguage, not an
object-language formula.  Functoids are introduced from the predefined
list below or via \texttt{def-functoid} (Section~\ref{sec:def-functoid}).
\end{remark}

\texttt{SET} is a primitive class constant.  The formula
\texttt{$a$ in SET} expresses ``$a$ is a set''; there is no separate
predicate form \texttt{SET($a$)}.
The key axiom distinguishing VNB from typed systems:
\[
  (a \in b) \;\Longrightarrow\; a \in \mathit{SET}
\]
Only sets can be members of anything: membership forces the \emph{left}
argument to be a set.  The right argument---the containing class---may be
a proper class.  ($3 \in \mathit{ORD}$ is well-formed even though the
class of all ordinals is a proper class.)

Function application is written $f(x_1, \ldots, x_n)$ in both cases:
for a function $f$ this abbreviates $\texttt{apply-function}(f, x_1, \ldots, x_n)$,
and for a functoid it abbreviates $\texttt{apply-functoid}(f, x_1, \ldots, x_n)$.

The subset relation
\[
  A \subseteq B \;\equiv\; \forall x.\; x \in A \Rightarrow x \in B
\]
is a primitive binary predicate written \texttt{A subset B}
(\texttt{subset-def}).  An immediate corollary of
\texttt{membership-implies-sethood} is that \emph{every class is a
subclass of $\mathit{SET}$} (axiom \texttt{subset-set}): only sets can
appear on the left of~$\in$.

The constructors below are \emph{total} functoids: each is defined for
\emph{any} class arguments and always produces a well-defined class.
Whether the result is a set depends on a separate \emph{sethood-closure}
condition on the arguments; those conditions are listed in
Section~\ref{sec:axioms} and Section~\ref{sec:funapp}.
Class comprehension $\{x \mid p\}$ is the sole exception: it is a
\emph{schema}, yielding one instance per syntactically correct formula
$p$, each taking $0$ arguments.

\begin{remark}
Totality is the VNB tradition.  \texttt{fun(BONGO, HARP)} is a
perfectly well-defined class---the class of all total functions from
\texttt{BONGO} to \texttt{HARP}---even when neither is a set.  The
axiom \texttt{fun-set-iff} then says that this class is itself a set
if and only if both \texttt{BONGO} and \texttt{HARP} are sets.
Definability and sethood are separate questions.
\end{remark}

\begin{remark}
A related discipline: when we add axioms about a constructor we state
them at the abstract level---how \texttt{fun}, \texttt{pair},
\texttt{tuples}, \texttt{list}, etc.\ \emph{behave}---never in terms
of how they might be concretely encoded as sets at the vanilla VNB
level.  So even if some particular encoding happened to make the list
\texttt{[2,3]} literally equal to the integer \texttt{23}, no theorem
of VNB+ may distinguish such coincidences.  The user should never
have to know whether ordered pairs are Kuratowski or Wiener, or
whether \texttt{nn} is built from von Neumann or Zermelo numerals.
\end{remark}

\begin{remark}[Variadic functoids]
Several of the constructors that follow---\texttt{union},
\texttt{intersection}, \texttt{cartesian}, and the type constructor
\texttt{fun}---are \emph{variadic functoids}.  The number of arguments
varies between invocations but is fixed when the expression is written:
any \emph{finite} number is allowed.

This is distinct from variadic \emph{functions} of type
$\mathit{FUN}(\mathit{TUPLES}(A), B)$, which take a single set-level
tuple as their one argument (see the remark on $\mathit{TUPLES}$ in
Section~\ref{sec:constructor-examples}), and from fixed-arity functoids
such as ring addition
$\mathit{ADD}(r) \in \mathit{FUN}(\mathit{CARTESIAN}(A(r), A(r)), A(r))$,
which have a single argument slot of fixed Cartesian type.
\end{remark}

\medskip\noindent The remaining constructors of this section are listed
below.  Each entry states the defining membership property and points to
the corresponding proof commands in
Section~\ref{sec:nary-proof-commands}.

\begin{itemize}[label = $\circ$, leftmargin=1cm]

\item \texttt{union($A_1$,\ldots,$A_n$)}
  $x \in C$ iff $x \in A_i$ for some $i$.
  Axiom \texttt{union-membership} states this IFF for $n=2$.
  At any arity, the kernel rule \texttt{pi-union-intro!} reduces a goal
  $x \in \text{union}(A_1,\ldots,A_n)$ to the single subgoal $x \in A_k$
  for a chosen index $k$; \texttt{pi-union-elim!} consumes a hypothesis
  $x \in \text{union}(A_1,\ldots,A_n)$ and case-splits it into $n$
  branches, each with $x \in A_i$ assumed.  See
  Section~\ref{sec:nary-proof-commands}.
  The macete \texttt{union-decompose} rewrites
  $x \in \text{union}(A_1,\ldots,A_n)$ to the disjunction
  $(x \in A_1) \lor \cdots \lor (x \in A_n)$ in one step.
  Result is a set when all $A_i \in \mathit{SET}$
  (\texttt{union-set-closure}; n-ary by induction).

\item \texttt{intersection($A_1$,\ldots,$A_n$)}
  $x \in C$ iff $x \in A_i$ for all $i$.
  Axiom \texttt{intersection-membership} states this IFF for $n=2$.
  At any arity, the kernel rule \texttt{pi-intersection-intro!} reduces
  a goal $x \in \text{intersection}(A_1,\ldots,A_n)$ to $n$ subgoals
  $x \in A_i$; \texttt{pi-intersection-elim!} extracts $x \in A_k$ from
  a hypothesis $x \in \text{intersection}(A_1,\ldots,A_n)$ for a chosen
  index $k$.  See Section~\ref{sec:nary-proof-commands}.
  The macete \texttt{intersection-decompose} rewrites
  $x \in \text{intersection}(A_1,\ldots,A_n)$ to the conjunction
  $(x \in A_1) \land \cdots \land (x \in A_n)$ in one step.
  Result is a set when any $A_i \in \mathit{SET}$
  (\texttt{intersection-set-closure}).

\item \texttt{complement-in($A$, $B$)}
  Relative complement $A \setminus B$: $x \in C$ iff $x \in A$ and
  $x \notin B$ (\texttt{complement-in-membership}).
  $B$ need not be a set; result is a set when $A \in \mathit{SET}$
  (\texttt{complement-in-set-closure}).

\item \texttt{cartesian($A_1$,\ldots,$A_n$)}
  $x \in C$ iff $x = [a_1,\ldots,a_n]$ with $a_i \in A_i$.
  Result is a set iff all $A_i \in \mathit{SET}$ (\texttt{cartesian-set-iff}).

\item \texttt{tuples($A$)}
  $x \in C$ iff $x = [a_1,\ldots,a_n]$ with all $a_i \in A$,
  for some $n \ge 0$.
  Result is a set when $A \in \mathit{SET}$ (\texttt{tuples-sethood}).

\item \texttt{fun($A$)}, \texttt{fun($A$, $B$)}
  The unary $\mathit{FUN}(A)$ is the class of all total functions with
  domain $A$ (no codomain restriction); proper class in general.  The
  binary $\mathit{FUN}(A,B)$ is the restriction
  $\{f \in \mathit{FUN}(A) : \forall x \in A.\;f(x) \in B\}$.
  $\mathit{FUN}(A,B) \in \mathit{SET}$ iff $A, B \in \mathit{SET}$
  (\texttt{fun-set-iff}); $\mathit{FUN}(A) \in \mathit{SET}$ never holds
  in general because the codomain is unconstrained.
  See Section~\ref{sec:funapp}.

  The strengthened \texttt{fun-domain-apply-def} now reads
  $f \in \mathit{FUN}(A) \Rightarrow ((f\,x) = (f\,x) \Leftrightarrow x \in A)$.
  Recall VNB equality is partial: $(t = t)$ is the definedness predicate
  for $t$, so the iff says \emph{$f(x)$ is defined exactly when $x \in A$}.
  Consequence: the witness $A$ is unique, justifying ``the'' domain of $f$.

\item \texttt{is-fun($f$)}
  The predicate ``$f$ is a function'': there exists a set $A$ with
  $f \in \mathit{FUN}(A)$.  Axiom \texttt{is-fun-def}:
  $\mathit{IS\text{-}FUN}(f) \Leftrightarrow \exists A \in \mathit{SET}.\;
   f \in \mathit{FUN}(A)$.

\item \texttt{dom($f$)}
  The \emph{domain} of $f$ as a class.  Characterised by membership
  (\texttt{dom-membership}):
  $x \in \mathit{DOM}(f) \Leftrightarrow x \in \mathit{SET}
   \land (f\,x) = (f\,x)$, i.e.\ $x$ is a set on which $f$ is defined.
  When $f \in \mathit{FUN}(A)$, $\mathit{DOM}(f)$ and $A$ agree on
  members (\texttt{dom-fun-membership}).  When additionally
  $A \in \mathit{SET}$, $\mathit{DOM}(f) \in \mathit{SET}$ and
  $\mathit{DOM}(f) = A$ as sets, by extensionality
  (\texttt{dom-of-fun}).

\item \texttt{res($f$, $B$)}
  The \emph{restriction} of $f$ to $B$.  Semantically
  $\mathit{RES}(f, B) = \lambda x \in B.\;f(x)$.  Total as a constructor
  (defined for all $f$ and $B$ -- bongo-style); typed by:
  \begin{itemize}[label=$-$, leftmargin=1cm]
  \item \texttt{res-typing}:
    $f \in \mathit{FUN}(A) \land B \subseteq A
    \Rightarrow \mathit{RES}(f, B) \in \mathit{FUN}(B)$.
  \item \texttt{res-apply}:
    $f \in \mathit{FUN}(A) \land B \subseteq A \land x \in B
    \Rightarrow (\mathit{RES}(f, B))(x) = f(x)$.
  \item \texttt{res-codomain} (derived): codomain inherits from
    $f \in \mathit{FUN}(A, C)$.
  \end{itemize}
  Outside the case $B \subseteq A$ the term exists but carries no
  $\mathit{FUN}$-typing guarantees.

\item \texttt{partial-fun($A$)}, \texttt{partial-fun($A$, $C$)}
  Partial functions on $A$.  Members are total functions whose domain
  is some subset of $A$:
  $\mathit{PARTIAL\text{-}FUN}(A) = \{f \mid \exists B \in \mathit{POWER}(A).
   \; f \in \mathit{FUN}(B)\}$
  (\texttt{partial-fun-membership}, generally a proper class because
  the codomain is unconstrained).  The binary form fixes a codomain:
  $\mathit{PARTIAL\text{-}FUN}(A, C) = \{f \mid \exists B \in \mathit{POWER}(A).
   \; f \in \mathit{FUN}(B, C)\}$
  (\texttt{partial-fun-binary-membership}).  Sethood
  (\texttt{partial-fun-binary-sethood}):
  $A, C \in \mathit{SET} \Rightarrow \mathit{PARTIAL\text{-}FUN}(A, C)
   \in \mathit{SET}$ (stated as axiom; provable once a class-union
  principle over a set-indexed family is in place).

\item \texttt{\{$x$ in $A$: $p$\}} / \texttt{\{$x$ in $A$ | $p$\}}
  Separation $\mathit{SEP}(x, A, p) = \{x \in A \mid p\}$: $x \in C$
  iff $x \in A$ and $p[x]$ holds.  Because $p$ is genuinely a schema
  (a formula, not a first-order term), the characterization is given
  by kernel rules rather than a single FOL axiom:
  \texttt{pi-sep-sethood!} ($\mathit{SEP}(x, A, p) \in \mathit{SET}$
  reduces to $A \in \mathit{SET}$);
  \texttt{pi-sep-mem-intro!} ($y \in \mathit{SEP}(x, A, p)$ splits
  into $y \in A$ and $p[x:=y]$);
  \texttt{pi-sep-mem-elim!} (assumption $y \in \mathit{SEP}(x, A, p)$
  yields both $y \in A$ and $p[x:=y]$ in context).
  Result is a set when $A \in \mathit{SET}$.

\item \texttt{\{$x$ | $p$\}}
  Class comprehension $\mathit{COMP}(x, p) = \{x \mid p\}$: $y \in C$
  iff $y \in \mathit{SET}$ and $p[x:=y]$ (members of any class are
  sets).  Same schema treatment as separation:
  \texttt{pi-comp-mem-intro!} and \texttt{pi-comp-mem-elim!}.
  Proper class in general unless provably bounded.

\item \texttt{power($A$)}, \texttt{power($A$, $B$)}
  Unary \texttt{power($A$)} is the power set: $x \in C$ iff
  $x \in \mathit{SET}$ and $x \subseteq A$ (\texttt{power-set}).
  Binary \texttt{power($A$, $B$)} is the function-space exponentiation
  $A^B$, axiomatized by \texttt{power-exp}:
  \[
     \mathit{POWER}(A, B) \;=\; \mathit{FUN}(B, A)
  \]
  i.e.\ functions $B \to A$.  This convention matches the arithmetic
  evaluator: $\texttt{power}(2, 3) = 2^3 = 8$ counts functions from
  a 3-set to a 2-set.  Sethood and membership inherit from
  \texttt{fun-set-iff} / \texttt{fun-codomain-iff}.

\end{itemize}

\begin{remark}
\texttt{union($A_1$, \ldots, $A_n$)} takes a finite listed union.
The \emph{big-union} $\bigcup A = \{x \mid \exists a \in A.\; x \in a\}$
of an arbitrary class $A$ of sets is a separate constructor
\texttt{UNION-SET}, currently pending (Section~\ref{sec:axioms}).
\end{remark}

\section{Term constructors}
\label{sec:choice-iota-lambda}

\emph{Term constructors} build syntactic objects whose intended
representation are objects in the VNB universe.

\begin{itemize}[label = $\circ$, leftmargin=1cm]
\item \texttt{[$e_1$, \ldots, $e_n$]} ordered $n$-tuple $(e_1, \ldots, e_n)$; always a set
\item \texttt{nth($k$, $e$)} $k$-th component of tuple $e$ (1-indexed); $k$ may be any term;
  \texttt{nth($k$, [$e_1$,\ldots,$e_n$])} reduces to $e_k$ via
  \texttt{cmd-nth-reduce} (Section~\ref{sec:commands})
\item \texttt{length($L$)} number of elements in tuple $L$; a natural number:
  \begin{itemize}[label=$-$, leftmargin=1cm]
  \item \texttt{length-of-empty}: $\texttt{length}([\,]) = 0$
  \item \texttt{length-in-nn}: $L \in \mathit{TUPLES}(\mathit{SET}) \Rightarrow
    \texttt{length}(L) \in \mathit{NN}$
  \item \texttt{nth-in-range}: $L \in \mathit{TUPLES}(A),\; i \in \mathit{NN},\;
    1 \le i \le \texttt{length}(L) \Rightarrow \texttt{nth}(i,L) \in A$
  \end{itemize}
  A full recursive characterisation requires a \texttt{prepend} constructor
  and the \texttt{tuples-induction} schema, both pending
  (Section~\ref{sec:pending-axioms}).
\item \texttt{EMPTY-SET} the empty set $\emptyset$; always a set; $= \texttt{make-set}([\,])$
\item \texttt{\{$a$\}} singleton $\{a\}$; always a set; $= \texttt{make-set}([$a$])$
\item \texttt{make-set($L$)} unordered finite set from tuple $L \in \mathit{TUPLES}(A)$,
  $A \in \mathit{SET}$; always a set:
  \begin{itemize}[label=$-$, leftmargin=1cm]
  \item \texttt{make-set-membership}: $x \in \texttt{make-set}(L)
    \Longleftrightarrow \exists\,i \in \mathit{NN}.\;
    1 \le i \le \texttt{length}(L) \land \texttt{nth}(i,L) = x$
    (the $1 \le i \le \texttt{length}(L)$ bounds matter:
    $\texttt{nth}$ out-of-range is unspecified)
  \item \texttt{make-set-sethood}: $L \in \mathit{TUPLES}(A),\; A \in \mathit{SET}
    \Rightarrow \texttt{make-set}(L) \in \mathit{SET}$
  \item \texttt{make-set-empty}: $\texttt{make-set}([\,]) = \emptyset$
  \end{itemize}
  The notation $\{a,b,c,\ldots\}$ is sugar for $\texttt{make-set}([a,b,c,\ldots])$;
  repetitions are harmless.
\item \texttt{choice($A$)} primitive choice operator: an element of nonempty class~$A$
  (undefined when $A$ is empty).  Axiom \texttt{choice-axiom}:
  \[
    (\exists x.\; x \in A) \;\Longrightarrow\; \texttt{choice}(A) \in A
  \]
  This is the Bourbaki-style global choice: a single canonical operator
  defined uniformly on all nonempty classes, not just sets.
\item \texttt{iota($x$, $p$)} definite description: the unique object satisfying~$p$;
  a binder.  Defined only when $\exists! x.\; p(x)$ is provable.  The
  characterization is the kernel rule \texttt{pi-iota-def!}, invoked
  with the iota term as argument; it posts two subgoals:
  \begin{enumerate}[label=(\arabic*),leftmargin=1.2cm]
  \item the existence-and-uniqueness obligation
    $\exists x.\;p \land (\forall y.\; p[x:=y] \Rightarrow x = y)$, and
  \item the original goal with the defining property
    $p[x := \texttt{iota}(x, p)]$ added as an assumption.
  \end{enumerate}
  A planned \texttt{define-object} form will wrap this pattern:
\begin{VNBsnip}
(define-object sqtwo (IOTA x (= (* x x) 2)))
; requires a proved unique-existence theorem; in this language,
; spelt out as (FORSOME x (AND (= (* x x) 2)
;                              (FORALL y (IMPLIES (= (* y y) 2) (= x y)))))
; (a FORSOME-UNIQUE macro abbreviation is planned but not yet defined)
; installs sqtwo satisfying: sqtwo * sqtwo = 2
\end{VNBsnip}
\item \texttt{lambda([$x_1$ in $A_1$, \ldots, $x_n$ in $A_n$], $\mathit{body}$)}
  anonymous multi-argument function; a binder.
  Each variable has its own independent domain, exactly as in quantification.
  All domains $A_i$ must belong to \texttt{SET}; the result is an element
  of $\mathit{FUN}$.  Two underlying representations:
  \begin{itemize}[label = $-$, leftmargin=1cm]
  \item \emph{record form} (\texttt{<functoid>}): produced by the parser
    when the lambda is given explicit binder domains
    \texttt{lambda([x in A], body)}.  Reduced by
    \texttt{cmd-functoid-beta} (short form \texttt{(beta)}).
  \item \emph{symbolic form} (\texttt{VNB-LAMBDA}): a raw S-expression
    binder \texttt{(VNB-LAMBDA $x$ $\mathit{body}$)} or
    \texttt{(VNB-LAMBDA (LIST $x_1\,\ldots\,x_n$) $\mathit{body}$)}.
    Used in axioms written directly as S-expressions
    (e.g.\ \texttt{RING-PROD}, \texttt{CC-RING}, \texttt{sum}).
    Typing rule \texttt{pi-lambda-type!} (single-binder shape):
    goal $\mathit{VNB\text{-}LAMBDA}(x, \mathit{body}) \in
    \mathit{FUN}(A, B)$ reduces to the subgoal
    $\forall x.\; x \in A \Rightarrow \mathit{body} \in B$.
    $\beta$-reduction \texttt{pi-lambda-beta!} (short form
    \texttt{(lam-b)}) rewrites
    $\bigl(\mathit{VNB\text{-}LAMBDA}(\overline x, \mathit{body})\bigr)
    (\overline v) = \mathit{body}[\overline x := \overline v]$
    (parallel substitution).
  \end{itemize}
  Cardinal-exponent typing schema (record-form, summarized; multi-binder
  derives from the single-binder rule plus uncurrying):
  \begin{multline*}
    \bigl(\forall x_i \in A_i.\;\mathit{body}(x_1,\ldots,x_n) \in B\bigr)
    \;\Longrightarrow\; \\
    \texttt{lambda}([x_1 \in A_1,\ldots,x_n \in A_n],\,\mathit{body})
      \in \mathit{FUN}(A_1 \times \cdots \times A_n,\,B)
  \end{multline*}

\item \texttt{lambdoid([$x_1$ in $A_1$, \ldots, $x_n$ in $A_n$], $\mathit{body}$)}
  like \texttt{lambda}, but the domains $A_i$ may be arbitrary classes
  (including proper classes such as \texttt{ORD}).  The result is tagged
  as a \emph{functoid} at the meta-level; it does not necessarily belong
  to any \texttt{FUN} class but still supports $\beta$-reduction via
  \texttt{(beta)}.
\end{itemize}

\begin{remark}[Functoids as meta-level objects]
Both \texttt{lambda} and \texttt{lambdoid} produce \emph{functoid records} at
the implementation level --- tagged Scheme structures distinct from plain
S-expressions and from VNB sets or classes.  This design follows the IMPS
principle that every structured entity carries its own recognisable tag; it
prevents the expression-traversal code from ever confusing a functoid with an
ordinary compound term.

Internally, the application
$\texttt{lambda}([x_1 \in A_1,\ldots], \mathit{body})(v_1,\ldots)$
is represented as
$(\texttt{apply-functoid}\;\langle\textit{record}\rangle\;v_1\;\cdots)$,
a mixed structure whose \texttt{car} is the symbol \texttt{apply-functoid}
and whose \texttt{cadr} is the record itself.  The surface syntax
$f(x)$ is used for \emph{all} application; the parser emits
\texttt{apply-functoid} only when $f$ is already tagged as a functoid.
\end{remark}

\begin{remark}
Tuples $[e_1,\ldots,e_n]$ are primitive here.  One could in principle reduce
them to iterated pairs (Kuratowski style), but we have no interest in doing so:
$[\ ]$ is perfectly well-defined as a term constructor of arity zero, and that
is all we need.
\end{remark}

\begin{remark}
The index $k$ in $\texttt{nth}(k, e)$ may be any term (not only a literal
integer).  The primitive inference \texttt{cmd-nth-reduce} does require a
literal integer to reduce $\texttt{nth}(k, [e_1,\ldots,e_n])$ to $e_k$, but
$\texttt{nth}(i, L)$ with a variable index $i$ is a perfectly valid term and
appears in \texttt{make-set-membership} and \texttt{nth-in-range}.
For example, $\texttt{nth}(3,\,[a,b,c,c,e])$ reduces to $c$.
\end{remark}

\begin{remark}[Indexing convention: lists are 1-indexed]
Tuples / lists in VNB are \emph{1-indexed}: \texttt{nth}$(i, L)$ is valid
for $1 \le i \le \texttt{length}(L)$, and $\texttt{nth}(1,\,[a, b, c]) = a$.
Lists and functions on initial segments of $\mathit{NN}$ are genuinely
different things, and we do \emph{not} reduce one to the other.  When an
interval constructor is added, the canonical domain for an $n$-element list
is $\mathit{INTERVAL}(1, n)$, \emph{not} $\mathit{INTERVAL}(0, n-1)$.
The 1-indexing convention matches IMPS and is built into every axiom
that mentions \texttt{nth} or \texttt{length} (\texttt{make-set-membership},
\texttt{nth-in-range}, \texttt{length-of-empty}).
\end{remark}

\begin{remark}
\texttt{choice}, \texttt{iota}, \texttt{lambda}, and \texttt{lambdoid} all
participate correctly in \texttt{free-vars}, \texttt{subst-free},
\texttt{alpha-equiv?}, and macete pattern-matching.  Both \texttt{lambda} and
\texttt{lambdoid} produce \emph{functoid records} --- Scheme tagged structures
separate from S-expressions --- so they are never accidentally mistaken for
ordinary terms by the expression machinery.
\end{remark}

\section{Function application}
\label{sec:funapp}

Function application is written as a compound term: \texttt{$f$($x$)} for
any object $f$ and any argument $x$.  Examples:
\begin{VNBsnip}
a + b               ; addition applied to a, b
NORM(E, v)          ; NORM of Banach space E applied to vector v
GEORGE(s)           ; syntactically valid, but may be undefined
\end{VNBsnip}

\texttt{fun}($A$, $B$) (with $A, B \in \mathit{SET}$) is a primitive set constructor: its members are
the total functions from $A$ to $B$.  Since $A, B \in \mathit{SET}$,
\texttt{fun}($A$,$B$) is itself a set (\texttt{fun-set-iff}), and each
$f \in \mathit{FUN}(A,B)$ is a set (\texttt{membership-implies-sethood}).

\begin{remark}
The soundness of $\mathit{FUN}(A,B)$ is confirmed by the standard
set-theoretic construction: a function $f : A \to B$ is a set of pairs
$(a, b) \in A \times B$ that is functional (each $a$ has at most one $b$)
and total (every $a \in A$ appears).  Since $A \times B$ is a set, so
is any subset of it, making $f$ a set.
\end{remark}


\begin{remark}
Category theory motivates a natural extension of $\mathit{FUN}$.  A
functor between large categories---such as the category of all groups or
all topological spaces---maps a proper class of objects to a proper class
of objects.  Since $\mathit{FUN}$ requires a set-valued domain, functors
fall outside $\mathit{FUN}$ entirely.

The planned treatment introduces \emph{operators} as a third kind of
entity alongside sets and classes.  There is no class $\mathit{OPER}(A,B)$
(no such collection can exist without violating the set/class dichotomy),
but there will be a predicate
\[
  \mathit{ISOPER}(A,\, B,\, f)
\]
asserting that $f$ behaves as a total function from $A$ to $B$ even when
$A$ or $B$ is a proper class.  An accompanying term constructor
$\texttt{apply-operator}(f, x)$ then gives the value of $f$ at $x$, for
any $x \in A$, under the hypothesis $\mathit{ISOPER}(A,B,f)$.

Operators are not members of anything---they do not appear on the left of
$\in$---so they are consistent with
\texttt{membership-implies-sethood}.  The three-level ontology (sets,
classes, operators) keeps the VNB foundation intact while accommodating
the full generality of categorical constructions.
\end{remark}

%
The base theory installs the following axioms.  The primary form of
\texttt{FUN} is the unary $\mathit{FUN}(A)$ -- all total functions whose
domain is $A$, with no codomain restriction; the binary $\mathit{FUN}(A,B)$
is the derived form $\{f \in \mathit{FUN}(A) : \forall x \in A.\;f(x) \in B\}$.

\begin{itemize}[label = $\circ$, leftmargin=1cm]
\item \texttt{fun-domain-apply-def}
  $f \in \mathit{FUN}(A) \land x \in A \;\Rightarrow\; f(x) = f(x)$
\item \texttt{fun-domain-extensionality}
  $f,g \in \mathit{FUN}(A) \land (\forall x \in A.\;f(x) = g(x)) \;\Rightarrow\; f = g$
\item \texttt{fun-codomain-iff}
  $f \in \mathit{FUN}(A,B) \;\Leftrightarrow\;
   (f \in \mathit{FUN}(A) \land \forall x \in A.\;f(x) \in B)$
\item \texttt{fun-apply-type}
  $f \in \mathit{FUN}(A,B) \land x \in A \;\Rightarrow\; f(x) \in B$
  (in \texttt{axioms.scm}; derivable from \texttt{fun-codomain-iff})
\item \texttt{fun-set-iff}
  $\mathit{FUN}(A,B) \in \mathit{SET} \;\Leftrightarrow\; (A \in \mathit{SET}) \land (B \in \mathit{SET})$
\item \texttt{fun-elements-are-sets}
  $f \in \mathit{FUN}(A,B) \land A \in \mathit{SET} \;\Rightarrow\; f \in \mathit{SET}$
\item \texttt{fun-domain-elements-are-sets}
  $f \in \mathit{FUN}(A) \land A \in \mathit{SET} \;\Rightarrow\; f \in \mathit{SET}$
\item \texttt{cartesian-set-iff}
  $A \times B \in \mathit{SET} \;\Leftrightarrow\; (A \in \mathit{SET}) \land (B \in \mathit{SET})$
\item \texttt{tuples-sethood}
  $A \in \mathit{SET} \;\Rightarrow\; \mathit{TUPLES}(A) \in \mathit{SET}$
\end{itemize}

\begin{remark}
The constructor $\mathit{TUPLES}(A)$ is the class of all finite sequences
$[a_1, \ldots, a_n]$ with each $a_i \in A$, for any $n \ge 0$ (including
the empty sequence $[\,]$).  It fills a gap that $\mathit{CARTESIAN}$ does
not cover: $\mathit{CARTESIAN}(A_1, \ldots, A_n)$ fixes the arity $n$ at
the time the expression is written, whereas $\mathit{TUPLES}(A)$ lets $n$
range over all natural numbers.

This supports \emph{variadic} functions whose single argument is a tuple
of any length: an element of $\mathit{FUN}(\mathit{TUPLES}(A), B)$ takes
one tuple and returns a value in $B$.  This is \emph{not} the same as a
fixed-arity $n$-ary function, which has type
$\mathit{FUN}(\mathit{CARTESIAN}(A_1, \ldots, A_n), B)$ and is bound to
exactly $n$ argument slots by its type.  The variadic flavour is the
right one when an operator must handle inputs of any length (e.g.\ a sum
over an indexed sequence); the fixed-arity flavour is the right one when
the arity is intrinsic to the operator (e.g.\ ring addition, which is
honestly binary).  Since $\mathit{TUPLES}(A) \in \mathit{SET}$ when
$A \in \mathit{SET}$, the variadic case is an ordinary instance of
$\mathit{FUN}(A', B)$ with $A' = \mathit{TUPLES}(A)$.
\end{remark}

\section{Constructor examples}
\label{sec:constructor-examples}

Quick reference for every class and term constructor.  Full descriptions
appear in Sections~\ref{sec:class-set-discipline}
and~\ref{sec:choice-iota-lambda}.

\paragraph{Functoids (class constructors).}  Each is defined for any class
arguments; the result is always a well-defined class.  Each row gives the
condition on $x$ for membership in the constructed class $C$, and the
sethood-closure condition on the arguments (when there is one).

\begin{center}
\footnotesize
\begin{tabular}{@{}p{4.5cm}p{4.8cm}p{3.0cm}@{}}
\toprule
Expression & $x \in C$ iff \ldots & $C \in \mathit{SET}$ when \\
\midrule
\texttt{union($A_1$,\ldots,$A_n$)} & $x \in A_i$ for some $i$ & all $A_i \in \mathit{SET}$ \\[2pt]
\texttt{intersection($A_1$,\ldots,$A_n$)} & $x \in A_i$ for all $i$ & any $A_i \in \mathit{SET}$ \\[2pt]
\texttt{complement-in($A$,$B$)} & $x \in A$ and $x \notin B$ & $A \in \mathit{SET}$ \\[2pt]
\texttt{cartesian($A_1$,\ldots,$A_n$)} & $x = [a_1,\ldots,a_n]$, $a_i \in A_i$ & all $A_i \in \mathit{SET}$ \\[2pt]
\texttt{tuples($A$)} & $x = [a_1,\ldots,a_n]$, $a_i \in A$, $n \ge 0$ & $A \in \mathit{SET}$ \\[2pt]
\texttt{fun($A$,$B$)} & $x$ is a total function from $A$ to $B$ & $A, B \in \mathit{SET}$ \\[2pt]
\texttt{\{$x$ in $A$: $p$\}} & $x \in A$ and $p(x)$ holds & $A \in \mathit{SET}$ \\[2pt]
\texttt{\{$x$ | $p$\}} & $p(x)$ holds & rarely \\[2pt]
\texttt{power($A$)} & $x \in \mathit{SET}$ and $x \subseteq A$ & $A \in \mathit{SET}$ \\
\bottomrule
\end{tabular}
\end{center}

\paragraph{Term constructors.}

\begin{center}
\footnotesize
\begin{tabular}{@{}p{5.2cm}p{7.1cm}@{}}
\toprule
Expression & Key property \\
\midrule
\texttt{EMPTY-SET} & empty set $\emptyset$; always a set; no members \\[2pt]
\texttt{\{$a$\}} & singleton $\{a\}$; always a set \\[2pt]
\texttt{\{$a$, $b$, \ldots\}} & finite set via \texttt{make-set}; repetitions harmless \\[2pt]
\texttt{[$e_1$, \ldots, $e_n$]} & ordered $n$-tuple; always a set \\[2pt]
\texttt{nth($k$, $e$)} & $k$-th component (1-indexed); reduces via \texttt{cmd-nth-reduce} \\[2pt]
\texttt{choice($A$)} & $(\exists x.\;x\in A) \Rightarrow \texttt{choice}(A) \in A$ \\[2pt]
\texttt{iota($x$, $p$)} & unique object satisfying $p$; requires $\exists!\,x.\;p(x)$ \\[2pt]
\texttt{lambda([$x_1$ in $A_1$,\ldots], $b$)} & multi-arg function; domains in \texttt{SET}; $\beta$-reduces via \texttt{(beta)} \\[2pt]
\texttt{lambdoid([$x_1$ in $A_1$,\ldots], $b$)} & like \texttt{lambda} but domains may be proper classes \\
\bottomrule
\end{tabular}
\end{center}

%% =========================================================

\section{Destructuring quantification}

The syntax extends the bounded-variable binding form
\texttt{[$x$ in $A$]} to allow the bound variable to be a
\emph{tuple pattern}:
\[
  \texttt{[[$n_1$, \ldots, $n_k$] in $A$]}
\]
Here $A$ is \emph{any class} --- a Cartesian product, a number system,
a separation class, a structure predicate, or any user-defined class.
There is no requirement that $A$ be a mathematical structure in the sense
of \texttt{def-structure}.

The binding introduces an anonymous implicit tuple variable~$t_0$, adds
$t_0 \in A$ to the assumptions, and binds each name $n_i$
to the $i$-th component of~$t_0$:
\[
  n_i \;\longmapsto\; \texttt{nth}(i,\; t_0)
\]
The variable~$t_0$ never appears to the user.  The scope of
$n_1, \ldots, n_k$ is \emph{lexical}: they have no meaning outside the
quantifier body.

\paragraph{Examples.}

Ranging over a Cartesian product (no structure involved):
\begin{VNBsnip}
forall([[x, y] in cartesian(RR, RR)], x + y = y + x)

; Expands to (internal form):
(FORALL t_0
  (IMPLIES (IN t_0 (CARTESIAN RR RR))
    (= (+ (NTH 1 t_0) (NTH 2 t_0))
       (+ (NTH 2 t_0) (NTH 1 t_0)))))
\end{VNBsnip}

Ranging over a named structure (the structure case is not special):
\begin{VNBsnip}
forall([[A, +, *, f] in NORMEDSPACE], v in A implies is-continuous(f, A, RR))

; Expands to (internal form):
(FORALL t_0
  (IMPLIES (IN t_0 NORMEDSPACE)
    (IMPLIES (IN v (NTH 1 t_0))
      (is-continuous (NTH 4 t_0) (NTH 1 t_0) RR))))
\end{VNBsnip}

Destructuring bindings compose freely with plain and bounded bindings
(Section~\ref{sec:multiquant}):
\begin{VNBsnip}
forall([[A, +, *, f] in NORMEDSPACE, v in A], is-continuous(f, A, RR))
\end{VNBsnip}


\section{Parser and printer}
\label{sec:parser-printer}

Four entry points provide string-level access to VNB formulas:

\begin{VNBsnip}
(make-wff-from-string str)  ; string -> <wff>  (validates via make-wff)
(parse-string str)          ; string -> raw S-expression (no validation)
(wff->string w)             ; <wff>  -> string
(pp w)                      ; display wff->string then newline
\end{VNBsnip}

\noindent
\texttt{parse-string} converts a string to a raw S-expression without
validation; it is used internally by the interactive commands and is
available for low-level \texttt{cmd-*} calls.  For normal use, pass
formula strings directly to \texttt{ai}, \texttt{bc}, \texttt{inst},
\texttt{cut}, \texttt{wk}, and \texttt{vnb-create-num-theorem}.
\texttt{make-wff-from-string} is the right choice when constructing a
formula to prove (for \texttt{sp}).

\paragraph{Operator precedence (high to low).}
\begin{center}
\small
\begin{tabular}{lll}
\toprule
Level & Operators & Associativity \\
\midrule
primary & atoms, \texttt{f(a,b)}, \texttt{[a,b]}, \texttt{(expr)} & --- \\
power & \texttt{\textasciicircum} & right \\
unary & \texttt{-x} & right \\
mul & \texttt{*}, \texttt{/} & left (n-ary \texttt{*}) \\
add & \texttt{+}, \texttt{-} & left (n-ary \texttt{+}) \\
cmp & \texttt{in}, \texttt{=}, \texttt{==}, \texttt{<=} & binary \\
and & \texttt{and} & left (n-ary) \\
or & \texttt{or} & left (n-ary) \\
iff & \texttt{iff} & left (n-ary) \\
implies & \texttt{implies} & right \\
\bottomrule
\end{tabular}
\end{center}

\paragraph{Division and negation sugar.}
\begin{VNBsnip}
x / y    =>  (* x (recip y))   ; division is syntactic sugar
-x       =>  (- x)             ; unary minus; lower prec than ^
-x ^ 2   =>  (- (power x 2))  ; i.e. -(x^2), not (-x)^2
\end{VNBsnip}

\paragraph{Quantifier syntax.}
Quantifiers use the binding-list form:
\begin{VNBsnip}
forall([x], body)              ; bare variable (all classes)
forall([x in A], body)         ; guarded variable
forall([[a, b, c] in A], body) ; tuple destructuring
forall([x in A, y, z in B], body)  ; mixed bindings
forsome([x in NN], body)       ; existential; same forms
\end{VNBsnip}

\noindent
The bracket syntax after the function name distinguishes quantifiers
from ordinary function calls.

\paragraph{Display conventions (\texttt{wff->string} / \texttt{pp}).}
\begin{itemize}[label=\textendash, leftmargin=1cm]
\item \begin{sloppypar}\textbf{Quantifiers} print in binding-list form:
  the primitive expanded form is reverse-engineered.
  \texttt{forall([x in nn, y], p(x, y))}\end{sloppypar}
\item \textbf{Logical connectives} \texttt{and}, \texttt{or}, \texttt{iff}
  print n-ary infix with surrounding parentheses:
  \texttt{((p and q) and r)}.
  \texttt{implies} prints binary infix: \texttt{(p implies q)}.
\item \textbf{Comparisons} \texttt{in}, \texttt{=}, \texttt{==},
  \texttt{<=} print binary infix: \texttt{(x in nn)}, \texttt{(x = y)}.
\item \textbf{Arithmetic} \texttt{+}, \texttt{*} print n-ary infix:
  \texttt{(x + y + z)}.  Binary \texttt{-} prints as \texttt{(x - y)};
  unary as \texttt{-x}.  Exponentiation as \texttt{(x \textasciicircum{} y)}.
\item \textbf{Lists} \texttt{[a, b, c]} in both input and output: the
  list syntax is the same in string form (internally stored as \texttt{LIST}).
\item \textbf{All other compound terms} use functional notation
  \texttt{f(x, y, z)}: \texttt{nth}, \texttt{cartesian}, \texttt{fun},
  \texttt{power} (unary set power), \texttt{not}, and user-defined heads.
  Nullary symbols print without parentheses.
\end{itemize}

\paragraph{Example round-trips.}
\begin{VNBsnip}
(pp (make-wff-from-string "forall([x in NN, y in NN], x + y = y + x)"))
; => forall([x in nn, y in nn], ((x + y) = (y + x)))

(pp (make-wff-from-string "not(P) implies falsity"))
; => (not(p) implies falsity)

(pp (make-wff-from-string "2 ^ 10 = 1024"))
; => ((2 ^ 10) = 1024)

(expression->string '(forall x (implies (in x nn)
                       (forall y (implies (in y nn) (= (+ x y) (+ y x)))))))
; => "forall([x in nn, y in nn], ((x + y) = (y + x)))"
\end{VNBsnip}

\noindent
All identifiers are case-folded to lowercase by MIT~Scheme, so
\texttt{NN} and \texttt{nn} are the same symbol; \texttt{FORALL} and
\texttt{forall} are also identical.  The printer always emits lowercase.

%% =========================================================

%% =========================================================
